% Source: Weinberg GR, physical PDF pp. 98--102
%         (printed pp. 73--77).
% Coverage: Section 3.3, Relation Between g_{\mu\nu} and
%           \Gamma^\lambda{}_{\mu\nu}; equations (3.3.1)--(3.3.10).
% Figures/tables/footnotes: none.
% Status: source-reviewed and compile-clean.
% Uncertainties: source description after Eq. (3.3.10) calls the proper-time
%                extremum ``usually a minimum''; retained verbatim.

\subsection{Relation Between
\texorpdfstring{\(g_{\mu\nu}\)}{g} and
\texorpdfstring{\(\Gamma^\lambda{}_{\mu\nu}\)}{Gamma}}\label{sec:3.3}

Our treatment of freely falling particles has shown that the field that
determines the gravitational force is the ``affine connection''
$\Gamma^\lambda{}_{\mu\nu}$, whereas the proper-time interval between two
events with a given infinitesimal coordinate separation is determined by the
``metric tensor'' $g_{\mu\nu}$. We now show that $g_{\mu\nu}$ is also the
gravitational potential; that is, its derivatives determine the field
$\Gamma^\lambda{}_{\mu\nu}$.

We first recall the formula for the metric tensor, Eq.~\eqref{eq:3.2.7}:
\[
  g_{\mu\nu}
  =\frac{\partial\xi^\rho}{\partial x^\mu}
  \frac{\partial\xi^\sigma}{\partial x^\nu}\eta_{\rho\sigma}.
\]
Differentiation with respect to $x^\lambda$ gives
\[
  \partial_\lambda g_{\mu\nu}
  =
  \frac{\partial^2\xi^\rho}{\partial x^\lambda\partial x^\mu}
  \frac{\partial\xi^\sigma}{\partial x^\nu}\eta_{\rho\sigma}
  +
  \frac{\partial\xi^\rho}{\partial x^\mu}
  \frac{\partial^2\xi^\sigma}{\partial x^\lambda\partial x^\nu}
  \eta_{\rho\sigma},
\]
and recalling \eqref{eq:3.2.11}, we have
\[
  \partial_\lambda g_{\mu\nu}
  =
  \Gamma^\kappa{}_{\lambda\mu}
  \frac{\partial\xi^\rho}{\partial x^\kappa}
  \frac{\partial\xi^\sigma}{\partial x^\nu}\eta_{\rho\sigma}
  +
  \Gamma^\kappa{}_{\lambda\nu}
  \frac{\partial\xi^\rho}{\partial x^\mu}
  \frac{\partial\xi^\sigma}{\partial x^\kappa}\eta_{\rho\sigma}.
\]
Using \eqref{eq:3.2.7} again, we find
\begin{equation}
  \partial_\lambda g_{\mu\nu}
  =\Gamma^\rho{}_{\lambda\mu}g_{\rho\nu}
  +\Gamma^\rho{}_{\lambda\nu}g_{\rho\mu}.
  \tag{3.3.1}\label{eq:3.3.1}
\end{equation}

Before solving for $\Gamma$, it is necessary to point out a subtlety in the
derivation of Eq.~\eqref{eq:3.3.1} that has been hidden by our too compact
notation. When we erect a locally inertial coordinate system $\xi^\rho(x)$,
we do so at a specific point $X$, and the coordinates that are locally
inertial at $X$ should be so labeled, as $\xi_X^\rho(x)$. Thus
Eqs.~\eqref{eq:3.2.7} and \eqref{eq:3.2.11} should properly be written as
\begin{equation}
  g_{\mu\nu}(X)=
  \left[
  \frac{\partial\xi_X^\rho(x)}{\partial x^\mu}
  \frac{\partial\xi_X^\sigma(x)}{\partial x^\nu}
  \eta_{\rho\sigma}\right]_{x=X},
  \tag{3.3.2}\label{eq:3.3.2}
\end{equation}
\begin{equation}
  \left[
  \frac{\partial^2\xi_X^\rho(x)}
  {\partial x^\mu\partial x^\nu}\right]_{x=X}
  =
  \Gamma^\lambda{}_{\mu\nu}(X)
  \left[\frac{\partial\xi_X^\rho(x)}
  {\partial x^\lambda}\right]_{x=X}.
  \tag{3.3.3}\label{eq:3.3.3}
\end{equation}
When we differentiate \eqref{eq:3.3.2} with respect to $X^\lambda$, we get
two kinds of terms. The first kind arises because we set $x=X$; these contain
just the second derivatives \eqref{eq:3.3.3} and can be easily calculated as
before. The second kind of term arises because $\xi_X^\rho(x)$ carries a
label $X$; these terms contain derivatives like
\begin{equation}
  \left[
  \frac{\partial^2\xi_X^\rho(x)}
  {\partial X^\lambda\partial x^\mu}\right]_{x=X},
  \tag{3.3.4}\label{eq:3.3.4}
\end{equation}
and do not seem to have anything to do with the metric or the affine
connection.

In order to deal with this second kind of term it is necessary to sharpen
somewhat our interpretation of what is meant by ``locally inertial'' in the
Principle of Equivalence. We shall see in Section~3.5 that first derivatives
of the metric tensor may be measured by comparing the rates of identical
clocks an infinitesimal spacetime distance apart. Hence we shall interpret
the Principle of Equivalence as meaning that \emph{the locally inertial
coordinates $\xi_X^\rho$ that we construct at a given point $X$ can be
chosen so that the first derivatives of the metric tensor vanish at $X$}. In
the coordinate system $\xi_X^\rho$ the metric tensor at a point $X'$ is
given by \eqref{eq:3.3.2} as
\[
  g^X_{\rho\sigma}(X')
  \equiv
  \left[
  \frac{\partial\xi_{X'}^\gamma(x)}
       {\partial\xi_X^\rho(x)}
  \frac{\partial\xi_{X'}^\delta(x)}
       {\partial\xi_X^\sigma(x)}
  \eta_{\gamma\delta}\right]_{x=X'},
\]
and our new interpretation of the Principle of Equivalence tells us that
this quantity is stationary in $X'$ at $X'=X$. In order to use this
information, we introduce an arbitrary ``laboratory'' coordinate system
$x^\mu$, and write
\begin{align*}
  g_{\mu\nu}(X')
  &\equiv
  \left[
  \frac{\partial\xi_{X'}^\gamma(x)}{\partial x^\mu}
  \frac{\partial\xi_{X'}^\delta(x)}{\partial x^\nu}
  \eta_{\gamma\delta}\right]_{x=X'}
  \\
  &=g^X_{\rho\sigma}(X')
  \left[
  \frac{\partial\xi_X^\rho(x)}{\partial x^\mu}
  \frac{\partial\xi_X^\sigma(x)}{\partial x^\nu}
  \right]_{x=X'}.
\end{align*}
Differentiating with respect to $X'^\lambda$ and setting $X'=X$ gives then
(because $g^X_{\rho\sigma}(X')$ is stationary)
\begin{align*}
  \frac{\partial g_{\mu\nu}(X)}{\partial X^\lambda}
  &=
  g^X_{\rho\sigma}(X)
  \left[
  \partial_\lambda\left\{
  \frac{\partial\xi_X^\rho(x)}{\partial x^\mu}
  \frac{\partial\xi_X^\sigma(x)}{\partial x^\nu}
  \right\}\right]_{x=X}
  \\
  &=
  \eta_{\rho\sigma}\left[
  \frac{\partial^2\xi_X^\rho(x)}
  {\partial x^\lambda\partial x^\mu}
  \frac{\partial\xi_X^\sigma(x)}{\partial x^\nu}
  +
  \frac{\partial\xi_X^\rho(x)}{\partial x^\mu}
  \frac{\partial^2\xi_X^\sigma(x)}
  {\partial x^\lambda\partial x^\nu}
  \right]_{x=X}.
\end{align*}
No derivatives like \eqref{eq:3.3.4} now appear, and we can use
\eqref{eq:3.3.2} and \eqref{eq:3.3.3} as before to show that
\[
  \frac{\partial g_{\mu\nu}(X)}{\partial X^\lambda}
  =
  \Gamma^\rho{}_{\lambda\mu}(X)g_{\rho\nu}(X)
  +\Gamma^\rho{}_{\lambda\nu}(X)g_{\rho\mu}(X),
\]
which is precisely Eq.~\eqref{eq:3.3.1}.

Now we return to our previous compact notation, and solve for the affine
connection. Add to Eq.~\eqref{eq:3.3.1} the same equation with $\mu$ and
$\lambda$ interchanged and subtract the same equation with $\nu$ and
$\lambda$ interchanged. We have then
\begin{align}
  \partial_\lambda g_{\mu\nu}
  +\partial_\mu g_{\lambda\nu}
  -\partial_\nu g_{\mu\lambda}
  ={}&
  g_{\kappa\nu}\Gamma^\kappa{}_{\lambda\mu}
  +g_{\kappa\mu}\Gamma^\kappa{}_{\lambda\nu}
  \notag\\
  &+g_{\kappa\nu}\Gamma^\kappa{}_{\mu\lambda}
  +g_{\kappa\lambda}\Gamma^\kappa{}_{\mu\nu}
  \notag\\
  &-g_{\kappa\lambda}\Gamma^\kappa{}_{\nu\mu}
  -g_{\kappa\mu}\Gamma^\kappa{}_{\nu\lambda}
  \notag\\
  ={}&2g_{\kappa\nu}\Gamma^\kappa{}_{\lambda\mu}.
  \tag{3.3.5}\label{eq:3.3.5}
\end{align}
(Recall that $\Gamma^\kappa{}_{\mu\nu}$ and $g_{\mu\nu}$ are symmetric under
interchange of $\mu$ and $\nu$.) Define a matrix $g^{\nu\sigma}$ as the
inverse of $g_{\kappa\nu}$, that is,
\begin{equation}
  g^{\nu\sigma}g_{\kappa\nu}=\delta^\sigma_\kappa,
  \tag{3.3.6}\label{eq:3.3.6}
\end{equation}
and multiply the above with $g^{\nu\sigma}$; this gives finally
\begin{equation}
  \Gamma^\sigma{}_{\lambda\mu}
  =\frac12g^{\nu\sigma}
  \left(\partial_\lambda g_{\mu\nu}
  +\partial_\mu g_{\lambda\nu}
  -\partial_\nu g_{\mu\lambda}\right).
  \tag{3.3.7}\label{eq:3.3.7}
\end{equation}

It should be noted that \eqref{eq:3.2.7} ensures that the metric tensor
\emph{does} have an inverse, given by
\begin{equation}
  g^{\nu\sigma}=g^{\sigma\nu}
  \equiv
  \eta^{\rho\kappa}
  \frac{\partial x^\nu}{\partial\xi^\rho}
  \frac{\partial x^\sigma}{\partial\xi^\kappa},
  \tag{3.3.8}\label{eq:3.3.8}
\end{equation}
for, using the familiar product rule
\[
  \frac{\partial x^\nu}{\partial\xi^\rho}
  \frac{\partial\xi^\gamma}{\partial x^\nu}
  =\delta^\gamma_\rho,
\]
we find
\begin{align*}
  g^{\nu\sigma}g_{\kappa\nu}
  &=
  \eta^{\rho\lambda}
  \frac{\partial x^\nu}{\partial\xi^\rho}
  \frac{\partial x^\sigma}{\partial\xi^\lambda}
  \eta_{\gamma\delta}
  \frac{\partial\xi^\gamma}{\partial x^\kappa}
  \frac{\partial\xi^\delta}{\partial x^\nu}
  \\
  &=
  \eta^{\rho\lambda}
  \frac{\partial x^\sigma}{\partial\xi^\lambda}
  \eta_{\gamma\rho}
  \frac{\partial\xi^\gamma}{\partial x^\kappa}
  =
  \frac{\partial x^\sigma}{\partial\xi^\gamma}
  \frac{\partial\xi^\gamma}{\partial x^\kappa}
  =\delta^\sigma_\kappa,
\end{align*}
as required by \eqref{eq:3.3.6}. The right-hand side of
Eq.~\eqref{eq:3.3.7} is often called a Christoffel symbol.

One important consequence of the relation between the affine connection and
the metric tensor is that the equation of motion of a freely falling particle
automatically maintains the form of the proper-time interval $\dd\tau$.
Using \eqref{eq:3.2.3} we may calculate that
\begin{align*}
  \frac{\dd}{\dd\tau}\left(
  g_{\mu\nu}\frac{\dd x^\mu}{\dd\tau}
  \frac{\dd x^\nu}{\dd\tau}\right)
  ={}&
  \partial_\lambda g_{\mu\nu}
  \frac{\dd x^\lambda}{\dd\tau}
  \frac{\dd x^\mu}{\dd\tau}
  \frac{\dd x^\nu}{\dd\tau}
  +g_{\mu\nu}\frac{\dd^2x^\mu}{\dd\tau^2}
  \frac{\dd x^\nu}{\dd\tau}
  +g_{\mu\nu}\frac{\dd x^\mu}{\dd\tau}
  \frac{\dd^2x^\nu}{\dd\tau^2}
  \\
  ={}&
  \left[\partial_\lambda g_{\mu\nu}
  -g_{\rho\nu}\Gamma^\rho{}_{\mu\lambda}
  -g_{\mu\rho}\Gamma^\rho{}_{\nu\lambda}\right]
  \frac{\dd x^\lambda}{\dd\tau}
  \frac{\dd x^\mu}{\dd\tau}
  \frac{\dd x^\nu}{\dd\tau},
\end{align*}
and \eqref{eq:3.3.5} tells us that this vanishes; that is,
\begin{equation}
  g_{\mu\nu}\frac{\dd x^\mu}{\dd\tau}
  \frac{\dd x^\nu}{\dd\tau}=-C,
  \tag{3.3.9}\label{eq:3.3.9}
\end{equation}
where $C$ is a constant of the motion. Hence once we choose initial
conditions such that $\dd\tau^2$ is given by \eqref{eq:3.2.6}, we have
$C=1$, and \eqref{eq:3.3.9} will ensure that \eqref{eq:3.2.6} continues to
hold along the particle's path. Similarly, for a massless particle the
initial conditions give $C=0$ (with $\tau$ replaced by some other parameter
$\sigma$), and the equations of motion will keep
$g_{\mu\nu}\dd x^\mu\dd x^\nu$ zero along the path.

An additional consequence of the relation \eqref{eq:3.3.5} is that we are
enabled to formulate the law of motion of freely falling bodies as a
variational principle. Let us introduce an arbitrary parameter $p$ to
describe the path and write the proper time elapsed when the particle falls
from point $A$ to $B$ as
\[
  T_{BA}=\int_A^B\frac{\dd\tau}{\dd p}\,\dd p
  =\int_A^B\left\{
  -g_{\mu\nu}\frac{\dd x^\mu}{\dd p}
  \frac{\dd x^\nu}{\dd p}\right\}^{1/2}\dd p.
\]
Now vary the path from $x^\mu(p)$ to $x^\mu(p)+\delta x^\mu(p)$, keeping
fixed the endpoints, that is, setting $\delta x^\mu=0$ at $p_A$ and $p_B$.
The change in $T_{BA}$ is
\[
  \begin{aligned}
  \delta T_{BA}
  ={}&\frac12\int_A^B
  \left\{-g_{\mu\nu}\frac{\dd x^\mu}{\dd p}
  \frac{\dd x^\nu}{\dd p}\right\}^{-1/2}
  \Biggl\{-\partial_\lambda g_{\mu\nu}\,\delta x^\lambda
  \frac{\dd x^\mu}{\dd p}\frac{\dd x^\nu}{\dd p}
  \\
  &\hspace{8em}
  -2g_{\mu\nu}\frac{\dd\delta x^\mu}{\dd p}
  \frac{\dd x^\nu}{\dd p}\Biggr\}\dd p.
  \end{aligned}
\]
The first factor within the integrand is just $\dd p/\dd\tau$, so the
integral can be rewritten as
\[
  \delta T_{BA}
  =-\int_A^B\left\{
  \frac12\partial_\lambda g_{\mu\nu}\,\delta x^\lambda
  \frac{\dd x^\mu}{\dd\tau}\frac{\dd x^\nu}{\dd\tau}
  +g_{\mu\nu}\frac{\dd\delta x^\mu}{\dd\tau}
  \frac{\dd x^\nu}{\dd\tau}\right\}\dd\tau.
\]
We now integrate by parts, neglecting the endpoint contributions because
$\delta x^\mu$ vanishes at $A$ and $B$. This gives
\[
  \delta T_{BA}
  =-\int_A^B\left\{
  \frac12\partial_\lambda g_{\mu\nu}
  \frac{\dd x^\mu}{\dd\tau}\frac{\dd x^\nu}{\dd\tau}
  -\partial_\sigma g_{\lambda\nu}
  \frac{\dd x^\sigma}{\dd\tau}\frac{\dd x^\nu}{\dd\tau}
  -g_{\lambda\nu}\frac{\dd^2x^\nu}{\dd\tau^2}
  \right\}\delta x^\lambda\,\dd\tau.
\]
Inserting Eq.~\eqref{eq:3.3.5} and recalling that
$\Gamma^\lambda{}_{\mu\nu}$ is symmetric in its lower indices, we find
\begin{equation}
  \delta T_{BA}
  =-\int_A^B\left\{
  \frac{\dd^2x^\nu}{\dd\tau^2}
  +\Gamma^\nu{}_{\mu\sigma}
  \frac{\dd x^\mu}{\dd\tau}\frac{\dd x^\sigma}{\dd\tau}
  \right\}g_{\lambda\nu}\,\delta x^\lambda\,\dd\tau.
  \tag{3.3.10}\label{eq:3.3.10}
\end{equation}
Hence the spacetime path taken by a particle that obeys the equations
\eqref{eq:3.2.3} for free fall will be such that the proper time elapsed is
an extremum (and usually a minimum), that is,
\[
  \delta T_{BA}=0.
\]
% VERIFY SOURCE: the source says ``usually a minimum.''  Timelike geodesics
% locally maximize proper time in the standard modern formulation.
We may therefore express the equations of motion \eqref{eq:3.2.3}
geometrically, by saying that a particle in free fall through the curved
spacetime called a gravitational field will move on the shortest (or
longest) possible path between two points, ``length'' being measured by the
proper time. Such paths are called \emph{geodesics}. For instance, we can
think of the sun as distorting spacetime just as a heavy weight distorts a
rubber sheet, and can consider a comet's path as being bent toward the sun to
keep the path as ``short'' as possible. However, this geometrical analogy is
an \emph{a posteriori} consequence of the equations of motion derived from
the equivalence principle, and plays no necessary role in our
considerations.
